\begin{lemma}[Claim 4.3(b2), missing-color clauses]
Let \(\mathcal R\) satisfy
\noderef{ExtensionBaseRestorationInvariant}, with initial coloring
\(\sigma_n\), output \(\varrho_1\), terminal support \(v_n\), and ladder
tree \(T_j\) containing \(v_n\).  Assume the following two explicit
preparations.
\begin{enumerate}
\item If \(v\in V(T_j)\) is not the support of an active slot, then
      \(\overline{\sigma}_n(v)=\overline{\varphi}_{j-1}(v)\).
\item For each nonterminal support \(w\in V(T_j)\), let
      \(i_0<\cdots<i_p\) be the exact support class, with exact fields
      \(\gamma_{i_a},\delta_{i_a}\).  Then
      \[
       \gamma_{i_{a+1}}=\delta_{i_a},
      \]
      the list
      \(\gamma_{i_0},\delta_{i_0},\ldots,\delta_{i_p}\) is injective,
      and
      \[
       \overline{\varphi}_{j-1}(w)
       =(\overline{\sigma}_n(w)\setminus\{\delta_{i_p}\})
          \cup\{\gamma_{i_0}\}.                               \tag{1}
      \]
\end{enumerate}
The first preparation and (1) are the two conclusions of
\noderef{ExtensionBaseClaimMissingColors}; the chain and injectivity are
the conclusions of \noderef{ExtensionBaseSupportClassColors}.  Then
\[
 \overline{\varrho}_1(v)=\overline{\varphi}_{j-1}(v)
       \quad(v\in V(T_j)\setminus\{v_n\}),                    \tag{2}
\]
and
\[
             \overline{\varrho}_1(v_n)
             =\overline{\sigma}_n(v_n).                       \tag{3}
\]
\end{lemma}
\begin{proof}
Write the trace colorings as
\(\sigma^0=\sigma_n,\ldots,\sigma^r=\varrho_1\).  First let
\(v\in V(T_j)\) be unsupported by every active slot.  At a restoration
step, \(T_j\) is an initial segment of \(T_t\), and the invariant says
that the switched component meets \(T_t\) only at its support \(v_t\).
Since \(v\ne v_t\), its vertex node does not belong to that component.
No edge incident with \(v\) can belong to it either, because such an
edge node would be adjacent to the vertex node in
\noderef{TwoColorIncidenceGraph}.  Thus every transition preserves all
edge colors incident with \(v\), and
\[
                 \overline{\varrho}_1(v)
                 =\overline{\sigma}_n(v)
                 =\overline{\varphi}_{j-1}(v),                \tag{4}
\]
where the second equality is preparation 1.

Now let \(w\ne v_n\) be an active support in \(T_j\), and take the exact
increasing class enumeration from preparation 2.  Its members all lie
in \(L^*=L\setminus L(v_n)\), and the exact range and descending-order
fields of \noderef{ExtensionBaseRestorationTrace} show that the steps
supported at \(w\) occur in the order
\[
                         i_p,i_{p-1},\ldots,i_0.               \tag{5}
\]
A step supported at a different vertex does not alter an edge incident
with \(w\), by the same exact-component-intersection argument used for
(4).  It therefore suffices to track the steps in (5).

Immediately before the step at \(i_a\), the invariant says that
\(\delta_{i_a}\) is missing at \(w\), and its path assertion supplies an
incident edge of color \(\gamma_{i_a}\); hence that gamma color is
present.  Expanding \noderef{TwoColorComponentSwap} at the root \(w\)
therefore gives
\[
 \overline{\sigma}^{\,\mathrm{after}}(w)
   =(\overline{\sigma}^{\,\mathrm{before}}(w)
        \setminus\{\delta_{i_a}\})\cup\{\gamma_{i_a}\}.        \tag{6}
\]
The chain
\(\gamma_{i_a}=\delta_{i_{a-1}}\) makes the new missing color exactly
the delta needed for the next step.  Injectivity of
\(\gamma_{i_0},\delta_{i_0},\ldots,\delta_{i_p}\) ensures that no later
deletion in (6) removes the final color \(\gamma_{i_0}\), and that no
earlier deletion or insertion changes an unrelated endpoint color.
Induction down (5) consequently telescopes (6) to
\[
 \overline{\varrho}_1(w)
   =(\overline{\sigma}_n(w)\setminus\{\delta_{i_p}\})
        \cup\{\gamma_{i_0}\}.                                 \tag{7}
\]
Preparation (1) identifies the right side of (7) with
\(\overline{\varphi}_{j-1}(w)\).

Every vertex in \(V(T_j)\setminus\{v_n\}\) is covered by (4) or (7),
which proves (2).  Finally, no restoration slot is supported at \(v_n\)
by the definition of \(L^*\).  Since \(v_n\in V(T_j)\subseteq V(T_t)\)
at every step, exact component intersection again says that the switched
component avoids its vertex node and all incident edge nodes.  Its
missing-color set is unchanged throughout the trace, proving (3) and
all missing-color clauses of (b2).
\end{proof}
